\section{Introduction}
\label{sec:intro}

\begin{figure}[t]
\centering
\fig{fig_hero}{\linewidth}
\Description{Two columns of CasePath screen crops. Left: the claim says the paediatrician suspects a link with the mould, and the next steps end with obtaining the medical confirmation. Right: the claim says it is not known whether the cough has anything to do with the mould, and the next steps begin with the question whether the claim alleges effects on anyone's health; the other five requests are the same in both.}
\caption{The running example in \casepath: one changed sentence in the claim turns the request for the medical confirmation into a question, and the other five requests are unchanged. Top: the claim reader's source view, with the quoted words highlighted. Bottom: ``Next steps'', written by the result checker, a model agent, from the checklist that code evaluated. The lease is attached to both claims, so it is present and not requested; Figure~\ref{fig:reconverge} shows the workflow without files. Unedited captures of the product at commit 6a72159c on two synthetic claims run with the tenancy defect workflow; cropped.}
\label{fig:hero}
\end{figure}

Sometimes a single sentence in a claim decides which documents a case handler should request. A tenant reports mould, and the insurer's rules require a medical confirmation if the mould has affected someone's health. One version of the claim says ``the paediatrician suspects a link with the mould''. Another says ``we do not know whether it has anything to do with the mould''. For the first, the handler should request the confirmation; for the second, hold it and ask. Every other request is the same (Figure~\ref{fig:hero}).

This is one instance of a class of problems: deciding which evidence to request when the decision depends on facts the file does not settle. We treat it as design-science research \cite{hevner2004design,gregor2013positioning}: we build an artefact for this class, state design principles with their boundary conditions, and evaluate the artefact in one confirmatory test.

A single prompt that applies the rules gives one answer per document and records no fact values that code could check. It does not distinguish a fact the file omits, one it denies and one it leaves open. Section~\ref{sec:rule} shows how a rule engine can keep them apart and what is lost when it does not.

\casepath, a product that prepares legal-protection claims for a case handler, separates reading from deciding. Agents read every fact several times, each value with a quotation. Code admits a value that decides the checklist only when every reading gives it and its quotation is in the file; an unmentioned fact takes the workflow's presumed value unless a reading quotes doubt. Code then executes the organisation's rules, written as a workflow, and holds exactly the documents whose need, given the admitted facts, depends on an open fact (Section~\ref{sec:rule} states the scope).

We tested the product at a frozen commit on sealed synthetic claims, under a pre-registered plan with \sthreeHypotheses{} hypotheses in a fixed order (Section~\ref{sec:test}). With the cost-efficient model, it made fewer errors than Rules in context, the same model applying the rules in a single prompt: on theft (H1, \sthreeHOneA{} against \sthreeHOneB{} wrong or missed request changes per pair of texts differing in one fact) and on tenancy (H2, \sthreeHTwoA{} against \sthreeHTwoB{} strict unsafe errors, requests and omissions the rules do not allow, per claim). It also made fewer tenancy errors than Rules in context on the frontier model (H3, \sthreeHThreeA{} against \sthreeHThreeB).

The theft rules had been reviewed for execution, with knowledge of the test conditions; with the original rules, Rules in context made \sthreeRicOrigThErrors{} errors per pair on the same model and \sthreeRicOrigFrontierThErrors{} on the frontier model. The tenancy references are produced by the engine \casepath implements, under a silence clause every design receives. As the plan predicted, most of the tenancy difference lies in the \sthreeClaimsDeniedOpen{} claims, which the test set was designed to include, whose defining event is denied or left open (\sthreeHTwoDeniedOpenA{} against \sthreeHTwoDeniedOpenB{} and \sthreeHThreeDeniedOpenB). On the other \sthreeClaimsOrdinary{} it did not differ significantly from the frontier model; against the same model they also differed (\sthreeHTwoOrdinaryA{} against \sthreeHTwoOrdinaryB, descriptive), which the plan did not predict. It held \sthreeCasePathTenCorrectHold\% of the documents the tenancy references hold, against \sthreeRicTenCorrectHold\% and \sthreeRicFrontierTenCorrectHold\%, and more that need not wait (\sthreeCasePathTenSpurHold{} against \sthreeRicTenSpurHold{} and \sthreeRicFrontierTenSpurHold{} per claim). H4 was not confirmed: on theft the product beat Model executes rules and ReAct, but not significantly Self-consistency, a majority vote over five samples of Rules in context; so H5 was not tested. The full product cost \sthreeCasePathTenRatio{} (tenancy) and \sthreeCasePathThRatio{} (theft) times the frontier single prompt; running only the components that can affect the checklist cost \sthreeCostChecklistShareTen\% and \sthreeCostChecklistShareTh\% of that, with no detectable difference in errors (the two runs shared some cached responses). The test does not show how the product does on real or long files, with wrong presumptions, or with a frontier model run in full.

We reached this design after two pre-registered studies of earlier versions failed, and we attributed each shortfall, post hoc, to a likely cause, such as faulty extracted rules or silence read as doubt (Section~\ref{sec:history}).

The article makes three contributions:
\begin{enumerate}[leftmargin=1.5em]
\item A decision rule that admits only agreeing, quoted readings of the facts that decide the checklist, and holds exactly the documents whose need depends on an open fact, with its properties and their scope (Section~\ref{sec:rule}).
\item \casepath, the artefact, and the design knowledge from building it over three studies: four design principles with their evidence and boundary conditions (Sections~\ref{sec:casepath}, \ref{sec:history} and~\ref{sec:discussion}).
\item A pre-registered confirmatory evaluation on sealed data against single-prompt and prior-work strategies and a frontier model (Section~\ref{sec:test}).
\end{enumerate}
As resources, we release the code, the research workflows, the synthetic data and the analysis \cite{casepath2026repo}, and the protocol of a planned study with claims experts on real files (Appendix~\ref{app:protocol}).
